% Generated draft: requires mathematical and source review.
\documentclass[11pt]{article}
\usepackage[margin=25mm]{geometry}
\usepackage{fontspec}
\setmainfont{Latin Modern Roman}
\usepackage{amsmath,amssymb,mathtools,unicode-math}
\setmathfont{Latin Modern Math}
\usepackage{enumitem,needspace,xurl,hyperref}
\setlength{\parindent}{0pt}
\setlength{\parskip}{5pt}
\setlength{\emergencystretch}{3em}
\newcommand{\N}{\mathbb{N}}
\newcommand{\Z}{\mathbb{Z}}
\newcommand{\Q}{\mathbb{Q}}
\newcommand{\R}{\mathbb{R}}
\newcommand{\C}{\mathbb{C}}
\newcommand{\F}{\mathbb{F}}
\newcommand{\A}{\mathbb{A}}
\newcommand{\PP}{\mathbb{P}}
\newcommand{\E}{\mathbb{E}}
\newcommand{\eps}{\varepsilon}
\newcommand{\dd}{\,\mathrm{d}}
\begin{document}
\section*{UnsolvedMath ID: 20000456}
\textbf{Source problem code:} catalogue entry
\subsection{Definitions and mathematical statement}
Poonen's AIM question fixes an elliptic curve $E\subset\PP^2$ over $\Q$ with $E(\Q)\cong\Z$, in particular rank one and no torsion. On the chosen affine coordinate chart define
\[
 S=\{P\in E(\Q):y(P)\text{ is defined and }y(P)=a^2+b^2\text{ for some }a,b\in\Q\}.
\]
It asks to describe $S$ and whether it is finite. More generally it asks which subsets can be realized as $\pi(X(\Q))\subset E(\Q)$ for maps $\pi:X\to E$. The precise affine/Weierstrass model and permitted source varieties are not specified. For example the affine conic-family equation $y(P)=a^2+b^2$ gives a higher-dimensional source with image $S$; a restriction to projective curves would be an additional hypothesis.

The report's rigidity statement concerns a smooth projective geometrically integral curve $C$ and a $\Q$-morphism $f:C\to E$. It claims that $f(C(\Q))$ is finite or a full coset
\[
 R+mE(\Q),\qquad R\in E(\Q),\quad m\geq1.
\]
The infinite case comes from a genus-one source with a rational point and a translated isogeny; a source without rational points has empty image. This is a statement about the entire rational image, not just containment in a progression.

Let $H_0=E(\Q)\cap E(\R)^0$, where the superscript denotes the identity real component. Combining the rigidity with a rank-one height bound attributed to Bhakta--Loughran--Rydin Myerson--Nakahara, the report claims that any infinite projective-curve image contained in $S$ must be disjoint from $H_0$. In particular, if $E(\Q)\subset E(\R)^0$, every such contained image is finite. Connectedness of $E(\R)$ is a sufficient condition for this inclusion, not a necessary one.

If the rational group meets both real components, then $H_0$ has index two. The report says the published sieve leaves possible odd cosets disjoint from $H_0$ uncontrolled. This does not construct such a coset inside $S$, prove $S$ infinite, or settle its central finiteness question. Nor does the curve-source result classify images from arbitrary higher-dimensional $X$.

The original remark states that the coordinate-band set $\{P\in E(\Q):1\leq y(P)\leq2\}$ is infinite and Diophantine. This is a distinct set, not the two-square locus. The exact model and real-coordinate assumptions and full references for both this remark and the height-sieve input are not supplied.
\subsection{Short English statement}
For a torsion-free rank-one elliptic curve, ask whether points whose $y$-coordinate is a sum of two rational squares form a finite set. The report restricts possible images of projective curves inside that locus to arithmetic-progressions, and excludes infinite such images meeting the identity real component. It leaves the original finiteness problem and general higher-dimensional images unresolved.
\subsection{Sources}
\begin{itemize}
\item Local registry: \url{data/unsolvedmath/problems.json}, permanent ID $20000456$, source code \texttt{catalogue entry}; corrected question, original remark and supplied research summary. Corrected-verified status is supplied metadata, not an independent source check here.
\item AIM workshop ``Rational and integral points on higher dimensional varieties'', Question 23, attributed to Poonen: \url{https://aimath.org/WWN/qptsurface2/qptsurface2.pdf}. Canonical location \texttt{aim-algebraic-number-theory-notes.json}, \texttt{notes[107]}; original literature field empty. No PDF was opened here.
\item Supplied height-bound attribution: Bhakta--Loughran--Rydin Myerson--Nakahara; no title, identifier, exact rate, coordinate assumptions or theorem statement supplied.
\item Supplied machine-generated research attempt 1: partial result, problem open at the run and currently partially solved; rigor metadata ``proved'', novelty confidence low. The non-peer-reviewed curve-source/sieve synthesis requires independent verification before citation.
\end{itemize}
\end{document}
